% Springer LNCS (Lecture Notes in Computer Science) template.
\documentclass{llncs}

\usepackage{amsmath}
\usepackage{booktabs}
\usepackage{graphicx}

\begin{document}

\title{A Springer LNCS Template}
\author{First Author\inst{1} \and Second Author\inst{2}}
\institute{\inst{1} Example University, Beijing, China
\email{first@example.edu}
\and
\inst{2} Example Institute, Shanghai, China
\email{second@example.edu}}

\maketitle

\begin{abstract}
LNCS papers use a single-column layout with a strict page limit. This
template demonstrates title/author blocks, sections, and a references
list in Springer style.
\end{abstract}

\section{Introduction}
LNCS is the standard format for many CS conferences and workshops.
Citations use the bracket style~\cite{spr_demo}.

\section{Method}
Present your approach. Math works as usual:
\begin{equation}
  \mathcal{L} = \sum_{i=1}^{N} \ell(f_\theta(x_i), y_i).
\end{equation}

\section{Results}
\begin{table}[h]
  \centering
  \caption{Example results.}
  \begin{tabular}{lcc}
    \toprule
    Method & Acc. & F1 \\
    \midrule
    Baseline & 81.2 & 0.79 \\
    Ours     & \textbf{88.1} & \textbf{0.87} \\
    \bottomrule
  \end{tabular}
\end{table}

\begin{thebibliography}{1}
\bibitem{spr_demo}
A. Author, B. Author: An example article. In: Proc. of a Conference, pp. 1--10 (2025)
\end{thebibliography}

\end{document}
